% =============================================================================
% information.tex
%
% Satu sumber kebenaran untuk data penulis, skripsi, institusi, pernyataan AI,
% pengesahan, dan komponen awal lain sesuai PPTA IPB 2026.
%
% Pemuatan: % =============================================================================
% information.tex
%
% Satu sumber kebenaran untuk data penulis, skripsi, institusi, pernyataan AI,
% pengesahan, dan komponen awal lain sesuai PPTA IPB 2026.
%
% Pemuatan: % =============================================================================
% information.tex
%
% Satu sumber kebenaran untuk data penulis, skripsi, institusi, pernyataan AI,
% pengesahan, dan komponen awal lain sesuai PPTA IPB 2026.
%
% Pemuatan: % =============================================================================
% information.tex
%
% Satu sumber kebenaran untuk data penulis, skripsi, institusi, pernyataan AI,
% pengesahan, dan komponen awal lain sesuai PPTA IPB 2026.
%
% Pemuatan: \input{config/information} dari src/main.tex.
% =============================================================================

% --- Data Penulis ---
% Nilai bawaan mengikuti teks placeholder pada templat resmi. Ganti seluruh
% placeholder ini dengan data sebenarnya sebelum naskah diserahkan.
\newcommand{\NamaPenulis}{Aldi Pramudya}
\newcommand{\NamaPenulisInggris}{Aldi Pramudya}
\newcommand{\NIM}{G6401231003}
\newcommand{\TempatLahir}{Kota Pontianak}
\newcommand{\TanggalLahir}{16 November 2026}
\newcommand{\EmailPenulis}{stezipramudya@apps.ipb.ac.id}

% --- Data Pembimbing ---
\newcommand{\NamaPembimbingSatu}{Dr. Medria Kusuma Dewi Hardhienata, S.Komp., Ph.D.}
\newcommand{\NIPPembimbingSatu}{}
\newcommand{\NamaPembimbingDua}{Dr. Adriani Sunuddin, S.Pi., M.Si.}
\newcommand{\NIPPembimbingDua}{}

% --- Data Penguji ---
\newcommand{\NamaPengujiSatu}{Nama Lengkap dan Gelar}
\newcommand{\NamaPengujiDua}{Nama Lengkap dan Gelar}

% --- Data Pengesahan ---
\newcommand{\NamaKetuaProgramStudi}{Dr. Sony Hartono Wijaya, S.Kom., M.Kom.}
\newcommand{\NIPKetuaProgramStudi}{19810809 200812 1 002}
% Pejabat kedua bersifat opsional. Isi true bila blok ini perlu ditampilkan.
\newcommand{\TampilkanPejabatKedua}{false}
\newcommand{\JabatanPejabatKedua}{Ketua Departemen/Wakil Dekan Bidang Akademik dan Kemahasiswaan}
\newcommand{\NamaPejabatKedua}{Nama Lengkap dan Gelar}
\newcommand{\NIPPejabatKedua}{}
% Biarkan kosong agar templat hanya menampilkan label dan ruang pengisian.
\newcommand{\TanggalUjian}{}
\newcommand{\TanggalLulus}{}

% --- Data Skripsi ---
\newcommand{\JudulSkripsi}{%
  Pemodelan Kesesuaian Habitat Hiu Paus (\textit{Rhincodon typus}) di Perairan Indonesia dengan Adaptasi Domain%
}
\newcommand{\JudulSkripsiInggris}{%
  Whale Shark (\textit{Rhincodon typus}) Habitat Suitability Modeling
  in Indonesian Waters
  with Domain Adaptation%
}
\newcommand{\KataKunci}{ditulis dalam bahasa Indonesia, disusun berdasarkan abjad, maksimum lima kata atau frasa}
\newcommand{\Keywords}{written in English, arranged alphabetically, maximum five words or phrases}

% Judul abstrak grafis maksimum enam kata, tidak termasuk kata depan/sambung.
\newcommand{\JudulAbstrakGrafis}{%
  Judul Abstrak Grafis Maksimum Enam Kata, Tidak Termasuk Kata Depan dan Kata Sambung}
\newcommand{\JudulAbstrakGrafisInggris}{%
  Graphical Abstract Title Maximum Six Words, Excluding Prepositions and Conjunctions}
% Isi dengan path berkas gambar untuk mengganti kotak placeholder.
\newcommand{\BerkasAbstrakGrafis}{}
\newcommand{\BerkasAbstrakGrafisInggris}{}

% --- Deklarasi penggunaan AI ---
% Nilai placeholder mempertahankan pilihan [menggunakan/tidak menggunakan]
% seperti pada DOTX. Ganti dengan true atau false setelah menentukan pilihan.
\newcommand{\MenggunakanAI}{true}
\newcommand{\NamaAlatAI}{berupa model bahasa besar (\textit{large language model}) Claude Opus 5.5, Gemini 4 Neon, Gemini 4 Hydrogen, GLM 5.3, dan DeepSeek 4.1 Flash}
\newcommand{\TujuanPenggunaanAI}{pembuatan draf dan rekayasa perangkat lunak penelitian}

% --- Gelar dan sumber dana ---
\newcommand{\GelarSarjana}{Sarjana Ilmu Komputer}
% Kosongkan bila penelitian tidak memperoleh pendanaan khusus.
\newcommand{\SumberDana}{}

% --- Data Institusi ---
\newcommand{\Departemen}{Ilmu Komputer}
\newcommand{\ProgramStudi}{Ilmu Komputer}
\newcommand{\ProgramStudiSingkat}{CS}
\newcommand{\Fakultas}{Sekolah Sains Data, Matematika, dan Informatika}
\newcommand{\Universitas}{Institut Pertanian Bogor}
\newcommand{\KotaUniversitas}{Bogor}
\newcommand{\Tahun}{2027}

% --- Data Waktu Pelaksanaan (untuk prakata) ---
\newcommand{\BulanMulaiPenelitian}{November}
\newcommand{\TahunMulaiPenelitian}{2026}
\newcommand{\BulanSelesaiPenelitian}{Maret}
\newcommand{\TahunSelesaiPenelitian}{2027}
\newcommand{\BulanTahunDokumen}{22 November 2026}
 dari src/main.tex.
% =============================================================================

% --- Data Penulis ---
% Nilai bawaan mengikuti teks placeholder pada templat resmi. Ganti seluruh
% placeholder ini dengan data sebenarnya sebelum naskah diserahkan.
\newcommand{\NamaPenulis}{Aldi Pramudya}
\newcommand{\NamaPenulisInggris}{Aldi Pramudya}
\newcommand{\NIM}{G6401231003}
\newcommand{\TempatLahir}{Kota Pontianak}
\newcommand{\TanggalLahir}{16 November 2026}
\newcommand{\EmailPenulis}{stezipramudya@apps.ipb.ac.id}

% --- Data Pembimbing ---
\newcommand{\NamaPembimbingSatu}{Dr. Medria Kusuma Dewi Hardhienata, S.Komp., Ph.D.}
\newcommand{\NIPPembimbingSatu}{}
\newcommand{\NamaPembimbingDua}{Dr. Adriani Sunuddin, S.Pi., M.Si.}
\newcommand{\NIPPembimbingDua}{}

% --- Data Penguji ---
\newcommand{\NamaPengujiSatu}{Nama Lengkap dan Gelar}
\newcommand{\NamaPengujiDua}{Nama Lengkap dan Gelar}

% --- Data Pengesahan ---
\newcommand{\NamaKetuaProgramStudi}{Dr. Sony Hartono Wijaya, S.Kom., M.Kom.}
\newcommand{\NIPKetuaProgramStudi}{19810809 200812 1 002}
% Pejabat kedua bersifat opsional. Isi true bila blok ini perlu ditampilkan.
\newcommand{\TampilkanPejabatKedua}{false}
\newcommand{\JabatanPejabatKedua}{Ketua Departemen/Wakil Dekan Bidang Akademik dan Kemahasiswaan}
\newcommand{\NamaPejabatKedua}{Nama Lengkap dan Gelar}
\newcommand{\NIPPejabatKedua}{}
% Biarkan kosong agar templat hanya menampilkan label dan ruang pengisian.
\newcommand{\TanggalUjian}{}
\newcommand{\TanggalLulus}{}

% --- Data Skripsi ---
\newcommand{\JudulSkripsi}{%
  Pemodelan Kesesuaian Habitat Hiu Paus (\textit{Rhincodon typus}) di Perairan Indonesia dengan Adaptasi Domain%
}
\newcommand{\JudulSkripsiInggris}{%
  Whale Shark (\textit{Rhincodon typus}) Habitat Suitability Modeling
  in Indonesian Waters
  with Domain Adaptation%
}
\newcommand{\KataKunci}{ditulis dalam bahasa Indonesia, disusun berdasarkan abjad, maksimum lima kata atau frasa}
\newcommand{\Keywords}{written in English, arranged alphabetically, maximum five words or phrases}

% Judul abstrak grafis maksimum enam kata, tidak termasuk kata depan/sambung.
\newcommand{\JudulAbstrakGrafis}{%
  Judul Abstrak Grafis Maksimum Enam Kata, Tidak Termasuk Kata Depan dan Kata Sambung}
\newcommand{\JudulAbstrakGrafisInggris}{%
  Graphical Abstract Title Maximum Six Words, Excluding Prepositions and Conjunctions}
% Isi dengan path berkas gambar untuk mengganti kotak placeholder.
\newcommand{\BerkasAbstrakGrafis}{}
\newcommand{\BerkasAbstrakGrafisInggris}{}

% --- Deklarasi penggunaan AI ---
% Nilai placeholder mempertahankan pilihan [menggunakan/tidak menggunakan]
% seperti pada DOTX. Ganti dengan true atau false setelah menentukan pilihan.
\newcommand{\MenggunakanAI}{true}
\newcommand{\NamaAlatAI}{berupa model bahasa besar (\textit{large language model}) Claude Opus 5.5, Gemini 4 Neon, Gemini 4 Hydrogen, GLM 5.3, dan DeepSeek 4.1 Flash}
\newcommand{\TujuanPenggunaanAI}{pembuatan draf dan rekayasa perangkat lunak penelitian}

% --- Gelar dan sumber dana ---
\newcommand{\GelarSarjana}{Sarjana Ilmu Komputer}
% Kosongkan bila penelitian tidak memperoleh pendanaan khusus.
\newcommand{\SumberDana}{}

% --- Data Institusi ---
\newcommand{\Departemen}{Ilmu Komputer}
\newcommand{\ProgramStudi}{Ilmu Komputer}
\newcommand{\ProgramStudiSingkat}{CS}
\newcommand{\Fakultas}{Sekolah Sains Data, Matematika, dan Informatika}
\newcommand{\Universitas}{Institut Pertanian Bogor}
\newcommand{\KotaUniversitas}{Bogor}
\newcommand{\Tahun}{2027}

% --- Data Waktu Pelaksanaan (untuk prakata) ---
\newcommand{\BulanMulaiPenelitian}{November}
\newcommand{\TahunMulaiPenelitian}{2026}
\newcommand{\BulanSelesaiPenelitian}{Maret}
\newcommand{\TahunSelesaiPenelitian}{2027}
\newcommand{\BulanTahunDokumen}{22 November 2026}
 dari src/main.tex.
% =============================================================================

% --- Data Penulis ---
% Nilai bawaan mengikuti teks placeholder pada templat resmi. Ganti seluruh
% placeholder ini dengan data sebenarnya sebelum naskah diserahkan.
\newcommand{\NamaPenulis}{Aldi Pramudya}
\newcommand{\NamaPenulisInggris}{Aldi Pramudya}
\newcommand{\NIM}{G6401231003}
\newcommand{\TempatLahir}{Kota Pontianak}
\newcommand{\TanggalLahir}{16 November 2026}
\newcommand{\EmailPenulis}{stezipramudya@apps.ipb.ac.id}

% --- Data Pembimbing ---
\newcommand{\NamaPembimbingSatu}{Dr. Medria Kusuma Dewi Hardhienata, S.Komp., Ph.D.}
\newcommand{\NIPPembimbingSatu}{}
\newcommand{\NamaPembimbingDua}{Dr. Adriani Sunuddin, S.Pi., M.Si.}
\newcommand{\NIPPembimbingDua}{}

% --- Data Penguji ---
\newcommand{\NamaPengujiSatu}{Nama Lengkap dan Gelar}
\newcommand{\NamaPengujiDua}{Nama Lengkap dan Gelar}

% --- Data Pengesahan ---
\newcommand{\NamaKetuaProgramStudi}{Dr. Sony Hartono Wijaya, S.Kom., M.Kom.}
\newcommand{\NIPKetuaProgramStudi}{19810809 200812 1 002}
% Pejabat kedua bersifat opsional. Isi true bila blok ini perlu ditampilkan.
\newcommand{\TampilkanPejabatKedua}{false}
\newcommand{\JabatanPejabatKedua}{Ketua Departemen/Wakil Dekan Bidang Akademik dan Kemahasiswaan}
\newcommand{\NamaPejabatKedua}{Nama Lengkap dan Gelar}
\newcommand{\NIPPejabatKedua}{}
% Biarkan kosong agar templat hanya menampilkan label dan ruang pengisian.
\newcommand{\TanggalUjian}{}
\newcommand{\TanggalLulus}{}

% --- Data Skripsi ---
\newcommand{\JudulSkripsi}{%
  Pemodelan Kesesuaian Habitat Hiu Paus (\textit{Rhincodon typus}) di Perairan Indonesia dengan Adaptasi Domain%
}
\newcommand{\JudulSkripsiInggris}{%
  Whale Shark (\textit{Rhincodon typus}) Habitat Suitability Modeling
  in Indonesian Waters
  with Domain Adaptation%
}
\newcommand{\KataKunci}{ditulis dalam bahasa Indonesia, disusun berdasarkan abjad, maksimum lima kata atau frasa}
\newcommand{\Keywords}{written in English, arranged alphabetically, maximum five words or phrases}

% Judul abstrak grafis maksimum enam kata, tidak termasuk kata depan/sambung.
\newcommand{\JudulAbstrakGrafis}{%
  Judul Abstrak Grafis Maksimum Enam Kata, Tidak Termasuk Kata Depan dan Kata Sambung}
\newcommand{\JudulAbstrakGrafisInggris}{%
  Graphical Abstract Title Maximum Six Words, Excluding Prepositions and Conjunctions}
% Isi dengan path berkas gambar untuk mengganti kotak placeholder.
\newcommand{\BerkasAbstrakGrafis}{}
\newcommand{\BerkasAbstrakGrafisInggris}{}

% --- Deklarasi penggunaan AI ---
% Nilai placeholder mempertahankan pilihan [menggunakan/tidak menggunakan]
% seperti pada DOTX. Ganti dengan true atau false setelah menentukan pilihan.
\newcommand{\MenggunakanAI}{true}
\newcommand{\NamaAlatAI}{berupa model bahasa besar (\textit{large language model}) Claude Opus 5.5, Gemini 4 Neon, Gemini 4 Hydrogen, GLM 5.3, dan DeepSeek 4.1 Flash}
\newcommand{\TujuanPenggunaanAI}{pembuatan draf dan rekayasa perangkat lunak penelitian}

% --- Gelar dan sumber dana ---
\newcommand{\GelarSarjana}{Sarjana Ilmu Komputer}
% Kosongkan bila penelitian tidak memperoleh pendanaan khusus.
\newcommand{\SumberDana}{}

% --- Data Institusi ---
\newcommand{\Departemen}{Ilmu Komputer}
\newcommand{\ProgramStudi}{Ilmu Komputer}
\newcommand{\ProgramStudiSingkat}{CS}
\newcommand{\Fakultas}{Sekolah Sains Data, Matematika, dan Informatika}
\newcommand{\Universitas}{Institut Pertanian Bogor}
\newcommand{\KotaUniversitas}{Bogor}
\newcommand{\Tahun}{2027}

% --- Data Waktu Pelaksanaan (untuk prakata) ---
\newcommand{\BulanMulaiPenelitian}{November}
\newcommand{\TahunMulaiPenelitian}{2026}
\newcommand{\BulanSelesaiPenelitian}{Maret}
\newcommand{\TahunSelesaiPenelitian}{2027}
\newcommand{\BulanTahunDokumen}{22 November 2026}
 dari src/main.tex.
% =============================================================================

% --- Data Penulis ---
% Nilai bawaan mengikuti teks placeholder pada templat resmi. Ganti seluruh
% placeholder ini dengan data sebenarnya sebelum naskah diserahkan.
\newcommand{\NamaPenulis}{Aldi Pramudya}
\newcommand{\NamaPenulisInggris}{Aldi Pramudya}
\newcommand{\NIM}{G6401231003}
\newcommand{\TempatLahir}{Kota Pontianak}
\newcommand{\TanggalLahir}{16 November 2026}
\newcommand{\EmailPenulis}{stezipramudya@apps.ipb.ac.id}

% --- Data Pembimbing ---
\newcommand{\NamaPembimbingSatu}{Dr. Medria Kusuma Dewi Hardhienata, S.Komp., Ph.D.}
\newcommand{\NIPPembimbingSatu}{}
\newcommand{\NamaPembimbingDua}{Dr. Adriani Sunuddin, S.Pi., M.Si.}
\newcommand{\NIPPembimbingDua}{}

% --- Data Penguji ---
\newcommand{\NamaPengujiSatu}{Nama Lengkap dan Gelar}
\newcommand{\NamaPengujiDua}{Nama Lengkap dan Gelar}

% --- Data Pengesahan ---
\newcommand{\NamaKetuaProgramStudi}{Dr. Sony Hartono Wijaya, S.Kom., M.Kom.}
\newcommand{\NIPKetuaProgramStudi}{19810809 200812 1 002}
% Pejabat kedua bersifat opsional. Isi true bila blok ini perlu ditampilkan.
\newcommand{\TampilkanPejabatKedua}{false}
\newcommand{\JabatanPejabatKedua}{Ketua Departemen/Wakil Dekan Bidang Akademik dan Kemahasiswaan}
\newcommand{\NamaPejabatKedua}{Nama Lengkap dan Gelar}
\newcommand{\NIPPejabatKedua}{}
% Biarkan kosong agar templat hanya menampilkan label dan ruang pengisian.
\newcommand{\TanggalUjian}{}
\newcommand{\TanggalLulus}{}

% --- Data Skripsi ---
\newcommand{\JudulSkripsi}{%
  Pemodelan Kesesuaian Habitat Hiu Paus (\textit{Rhincodon typus}) di Perairan Indonesia dengan Adaptasi Domain%
}
\newcommand{\JudulSkripsiInggris}{%
  Whale Shark (\textit{Rhincodon typus}) Habitat Suitability Modeling
  in Indonesian Waters
  with Domain Adaptation%
}
\newcommand{\KataKunci}{ditulis dalam bahasa Indonesia, disusun berdasarkan abjad, maksimum lima kata atau frasa}
\newcommand{\Keywords}{written in English, arranged alphabetically, maximum five words or phrases}

% Judul abstrak grafis maksimum enam kata, tidak termasuk kata depan/sambung.
\newcommand{\JudulAbstrakGrafis}{%
  Judul Abstrak Grafis Maksimum Enam Kata, Tidak Termasuk Kata Depan dan Kata Sambung}
\newcommand{\JudulAbstrakGrafisInggris}{%
  Graphical Abstract Title Maximum Six Words, Excluding Prepositions and Conjunctions}
% Isi dengan path berkas gambar untuk mengganti kotak placeholder.
\newcommand{\BerkasAbstrakGrafis}{}
\newcommand{\BerkasAbstrakGrafisInggris}{}

% --- Deklarasi penggunaan AI ---
% Nilai placeholder mempertahankan pilihan [menggunakan/tidak menggunakan]
% seperti pada DOTX. Ganti dengan true atau false setelah menentukan pilihan.
\newcommand{\MenggunakanAI}{true}
\newcommand{\NamaAlatAI}{berupa model bahasa besar (\textit{large language model}) Claude Opus 5.5, Gemini 4 Neon, Gemini 4 Hydrogen, GLM 5.3, dan DeepSeek 4.1 Flash}
\newcommand{\TujuanPenggunaanAI}{pembuatan draf dan rekayasa perangkat lunak penelitian}

% --- Gelar dan sumber dana ---
\newcommand{\GelarSarjana}{Sarjana Ilmu Komputer}
% Kosongkan bila penelitian tidak memperoleh pendanaan khusus.
\newcommand{\SumberDana}{}

% --- Data Institusi ---
\newcommand{\Departemen}{Ilmu Komputer}
\newcommand{\ProgramStudi}{Ilmu Komputer}
\newcommand{\ProgramStudiSingkat}{CS}
\newcommand{\Fakultas}{Sekolah Sains Data, Matematika, dan Informatika}
\newcommand{\Universitas}{Institut Pertanian Bogor}
\newcommand{\KotaUniversitas}{Bogor}
\newcommand{\Tahun}{2027}

% --- Data Waktu Pelaksanaan (untuk prakata) ---
\newcommand{\BulanMulaiPenelitian}{November}
\newcommand{\TahunMulaiPenelitian}{2026}
\newcommand{\BulanSelesaiPenelitian}{Maret}
\newcommand{\TahunSelesaiPenelitian}{2027}
\newcommand{\BulanTahunDokumen}{22 November 2026}
